% Generated draft: requires mathematical and source review.
\documentclass[11pt]{article}
\usepackage[margin=25mm]{geometry}
\usepackage{fontspec}
\setmainfont{Latin Modern Roman}
\usepackage{amsmath,amssymb,mathtools,unicode-math}
\setmathfont{Latin Modern Math}
\usepackage{enumitem,needspace,xurl,hyperref}
\setlength{\parindent}{0pt}
\setlength{\parskip}{5pt}
\setlength{\emergencystretch}{3em}
\newcommand{\N}{\mathbb{N}}
\newcommand{\Z}{\mathbb{Z}}
\newcommand{\Q}{\mathbb{Q}}
\newcommand{\R}{\mathbb{R}}
\newcommand{\C}{\mathbb{C}}
\newcommand{\F}{\mathbb{F}}
\newcommand{\A}{\mathbb{A}}
\newcommand{\PP}{\mathbb{P}}
\newcommand{\E}{\mathbb{E}}
\newcommand{\eps}{\varepsilon}
\newcommand{\dd}{\,\mathrm{d}}
\begin{document}
\section*{UnsolvedMath ID: 9700011}
\textbf{Source problem code:} catalogue entry
\subsection{Definitions and mathematical statement}
Let \(\Pi\) be a homogeneous planar Poisson process with fixed intensity \(\lambda>0\), and let \(\mathcal G(\Pi)\) be the intended proximity graph, with straight edges weighted by Euclidean length. The graph rule is not specified; Delaunay, Gabriel, and relative-neighborhood graphs are examples mentioned in the supplied research summary, not interchangeable definitions.

One precise endpoint convention is the two-point Palm model obtained by adjoining \(o=(0,0)\) and \(z_r=(r,0)\) before constructing the graph. For a graph path \(\gamma=(x_0,\ldots,x_m)\) from \(o\) to \(z_r\), set
\[
\ell(\gamma)=\sum_{i=1}^m\|x_i-x_{i-1}\|,
\qquad D_r=\inf_{\gamma:o\to z_r}\ell(\gamma).
\]
For an attained shortest route \(\Gamma_r\), with a specified tie rule if needed, a perpendicular-deviation convention is
\[
T_r=\sup\{|y|:(x,y)\in\Gamma_r\}.
\]
Distance to the finite endpoint segment, rather than to its supporting line, is another convention and is not silently identified with this one.

The two requested asymptotic tasks are rigorous orders of transversal deviation and of \(\operatorname{Var}(D_r)\) as \(r\to\infty\). If the first is a mean-order question, one would seek deterministic scales \(a(r),b(r)\) and positive constants with
\[
\mathbb E T_r\asymp a(r),\qquad\operatorname{Var}(D_r)\asymp b(r),
\]
where \(\asymp\) means two-sided constant-factor bounds for all sufficiently large \(r\). A typical/high-probability order for \(T_r\) is a different target. No exponents or conjectural scales are supplied, and finite second moments/route existence must be part of the selected model.

The source does not specify graph rule, endpoint sampling, length metric, tie rule, or the probabilistic meaning of order for \(T_r\). Deterministic points are almost surely not original Poisson points, so adding anchors cannot be silently replaced by nearest-point routing. The supplied research summary reports no located published rigorous orders in the intended models, while citing related work on different random-geometric metrics; those results are not transferred here.
\subsection{Short English statement}
Find rigorous large-distance scales for shortest-route transverse fluctuations and route-length variance in the specified planar Poisson proximity graph. Fix the graph and endpoint/deviation conventions before asserting an order or exponent.
\subsection{Sources}
\begin{itemize}
\item Local UnsolvedMath registry, permanent ID \(9700011\), exact source code \texttt{catalogue entry}, ``Shortest routes in random proximity networks.'' Proposer David Aldous.
\item Aldous, \emph{Open problems}, page identified as \texttt{random\_proximity.html}; index \url{https://www.stat.berkeley.edu/~aldous/Research/OP/index.html}. Recorded access 2026-07-29; HTML text with linked-primary-document extraction.
\item Supplied OPEN-TRIAGE summary mentions Aldous's ``Scale-invariant random spatial networks'' and ``Shortest routes through networks,'' and related Bhamidi--van der Hofstad work on different random-geometric models, without complete matching references or asymptotic theorems.
\end{itemize}
\end{document}
